\section{Interface to Control Algorithms}

\subsection{Minimal Solver API}

For controller, simulator, and model-based RL integration, the learned model should eventually expose a minimal solver-style interface:
\begin{align}
\mathrm{reset}(\hat{\mathbf{x}}_0, \mathbf{c}_0) &\rightarrow \text{internal state}, \\
\mathrm{step}(\mathbf{u}_k, \Delta t, \mathbf{c}_k) &\rightarrow (\hat{\mathbf{x}}_{k+1}, \mathbf{i}_{k+1}), \\
\mathrm{rollout}(\mathbf{u}_{k:k+H-1}, \mathbf{c}_{k:k+H-1}) &\rightarrow \hat{\mathbf{x}}_{k+1:k+H}.
\end{align}

Here, $\mathbf{i}_{k+1}$ denotes diagnostic information returned together with the predicted next state.

\subsection{Diagnostic Information}

The diagnostic payload should include at least:
\begin{itemize}
  \item sensor freshness indicators,
  \item quality-context values such as elapsed time since the last DVL correction,
  \item predictive uncertainty or confidence surrogates when available,
  \item non-finite-event flags and fallback triggers,
  \item timing information for control-loop accounting.
\end{itemize}

This is important because a solver that predicts a state value without exposing reliability information is difficult to use safely inside controller or RL loops.

\subsection{Compatibility With Controller and RL Loops}

In a controller loop, the model is recursively queried as
\begin{equation}
\hat{\mathbf{x}}_{k+1}
=
f_{\theta}\!\left(\hat{\mathbf{x}}_{k}, \mathbf{u}_{k}, \mathbf{c}_{k}\right),
\end{equation}
and the resulting state is used for cost evaluation, policy updates, or search over future actions.

In model-based RL, the same interface supports:
\begin{itemize}
  \item imagined rollouts,
  \item short-horizon planning,
  \item synthetic trajectory generation for policy improvement.
\end{itemize}

Therefore, the main engineering requirement is not a visually impressive trajectory plot, but a numerically stable and well-instrumented transition step.

\subsection{Current Repository Status}

The repository does not yet expose a fully unified production-ready \texttt{step(u\_t, dt)} interface.
However, the present data contract, the single-step transition branch, and the offline transition-solver wrapper already define the mathematical form such an interface should follow.

In other words, the project is no longer only discussing a solver in the abstract:
it already has the ingredients for offline replay-style solver validation,
but it has not yet completed the final interface unification needed for controller deployment.

This makes the current phase a preparation stage for solver-style deployment rather than a terminal modeling stage.
